% Abstract (about 290 words) - paste into the ABSTRACT page of the template
Learners who study from online videos often need help to understand difficult concepts and to decide what
to study next, and AI assistants that provide this help depend on network connectivity and computing
resources that vary, especially for learners outside cities. This project proposes a QoE-aware agentic
framework for adaptive digital learning assistance that keeps answers correct, grounded and timely as
these conditions change. The framework answers questions from a lesson's notes and the transcript of the
chosen video using retrieval-augmented generation with corrective retrieval, generates each answer either
with a small language model on the learner's laptop or with a large model in the cloud according to a
runtime policy, measures every session as quality of experience (QoE) and resource efficiency, and selects
the next session's policy from similar past sessions. A proactive coach sets concept-level goals from a
pre-test, suggests the next step after every answer and reflects after a parallel post-test; speech input
and output support accessibility. The approach was implemented for five NCERT Class X Biology chapters and
evaluated without learners, using real answers and timings on 75 keyed questions and 40 learner-style and
off-topic questions, measured and emulated network conditions, and 10,560 replayed sessions. The cloud model
answered all keyed items correctly against 54\% for the edge model; corrective retrieval raised the
faithfulness of edge answers from 0.70 to 0.80; and the revised policy selector (QAPS) reached a session
utility of 0.812, above every fixed policy (best 0.794) and close to the hindsight optimum (0.818), with no
failed cloud calls and 44\% less energy than always using the cloud. Learning gain was not measured and
requires a user study. The project contributes to SDG 4 (Quality Education), targets 4.5 and 4.a, by
making grounded learning support usable on ordinary laptops and unreliable networks; this is an expected
benefit, not a demonstrated impact on learning.

Keywords: Quality of Experience; retrieval-augmented generation; edge--cloud inference; adaptive policy
selection; proactive tutoring
